% ACL BioNLP 2026 Workshop Submission
% 8-page limit (excluding references)
% Template: ACL 2026

\documentclass[11pt]{article}
\usepackage[hyperref]{acl}
\usepackage{times}
\usepackage{latexsym}
\usepackage{booktabs}
\usepackage{amsmath}
\usepackage{graphicx}

\title{Evaluation of RAG Knowledge Bases for Smoking Cessation Counseling: A Data Leakage Warning}

\author{Research Team \\
  Department of Computer Science \\
  University Name \\
  \texttt{email@university.edu}}

\begin{document}
\maketitle

\begin{abstract}
We present a rigorous evaluation of Retrieval-Augmented Generation (RAG) approaches for smoking cessation counseling, discovering severe test set contamination that completely reversed initial conclusions. After creating an independent test set validated for zero data leakage (n=15 questions, <50\% similarity threshold), we found that: (1) baseline LLM outperformed all RAG variants (ROUGE-L: 0.231 vs 0.224 best RAG), (2) AI-generated knowledge bases matched human-curated ones (102\% relative performance, p=0.184), and (3) 47\% of the original test set was severely contaminated ($\geq$90\% similarity to training data). These findings demonstrate the critical importance of test set independence verification and suggest RAG benefit is domain-dependent.
\end{abstract}

\section{Introduction}

Retrieval-Augmented Generation (RAG) has emerged as a promising approach to enhance LLM performance by incorporating external knowledge bases \cite{lewis2020retrieval}. However, rigorous evaluation requires test sets with zero contamination from training data.

We evaluate RAG for smoking cessation counseling—a well-documented domain with established clinical guidelines \cite{fiore2008treating}. While seemingly ideal for RAG, this comprehensive documentation may paradoxically reduce benefit if LLMs have already learned this knowledge during pre-training.

\textbf{Key Contributions:}
\begin{itemize}
    \item \textbf{Methodological warning:} Discovery of severe test set contamination (47\% of questions $\geq$90\% similar to training data) that completely reversed conclusions
    \item \textbf{Rigorous evaluation:} Independent test set with <50\% similarity validation, 5 runs, comprehensive statistical testing
    \item \textbf{Surprising finding:} Baseline outperforms RAG for well-documented domains, proposing a domain saturation hypothesis
    \item \textbf{Cost-effective alternative:} AI-generated knowledge bases match human-curated performance
\end{itemize}

\section{Related Work}

\textbf{RAG Systems.} RAG enhances LLMs by retrieving relevant documents at inference time \cite{lewis2020retrieval}. Recent work explores dense retrieval, hybrid search, and domain adaptation.

\textbf{Data Leakage in NLP.} Test-train overlap has been identified in QA datasets \cite{lewis2021questionanswering}, benchmark contamination in language models \cite{dodge2021documenting}, and memorization in neural networks. Our work demonstrates how leakage can reverse evaluation conclusions.

\textbf{Health NLP.} Prior work applies NLP to clinical text, patient education, and health question answering. We extend this to smoking cessation counseling with rigorous evaluation methodology.

\section{Data Leakage Discovery}

During preliminary analysis, we discovered significant overlap between the original test set and training data using fuzzy string matching (Python's \texttt{SequenceMatcher}):

\begin{table}[h]
\centering
\small
\begin{tabular}{lcc}
\toprule
\textbf{Severity} & \textbf{Count} & \textbf{\%} \\
\midrule
Severe ($\geq$90\%) & 7 & 47\% \\
Moderate (70-90\%) & 6 & 40\% \\
Clean (<70\%) & 2 & 13\% \\
\textbf{Exact matches (100\%)} & \textbf{4} & \textbf{27\%} \\
\bottomrule
\end{tabular}
\caption{Original test set contamination analysis}
\label{tab:leakage}
\end{table}

This contamination fundamentally invalidated the original evaluation, as test questions were essentially memorized during training. \textbf{Critically, the contaminated test set showed a completely different ranking}: Raw $>$ Human $>$ AI (wrong), versus the true ranking: Baseline $>$ AI $>$ Human $>$ Raw (correct).

\section{Methodology}

\subsection{Independent Test Set Creation}

We created a new test set with rigorous independence validation:

\begin{itemize}
    \item \textbf{Similarity threshold:} <50\% fuzzy string match to any training question
    \item \textbf{Sample size:} 15 questions covering diverse topics (benefits, cravings, NRT, withdrawal, relapse, mechanisms)
    \item \textbf{Validation:} Domain expert-reviewed reference answers
    \item \textbf{Grammatical diversity:} Distinct phrasings while maintaining topic coverage
\end{itemize}

\subsection{RAG Approaches}

We compared four approaches:

\begin{enumerate}
    \item \textbf{Baseline:} No retrieval, direct LLM responses
    \item \textbf{Raw Sources:} Retrieval from unprocessed CDC guidelines and clinical papers
    \item \textbf{AI-Generated:} Retrieval from GPT-4-generated Q\&A pairs
    \item \textbf{Human-Curated:} Retrieval from expert-curated Q\&A database
\end{enumerate}

\subsection{Evaluation Protocol}

\begin{itemize}
    \item \textbf{Model:} GPT-4o-mini (temperature=0.3, max\_tokens=200)
    \item \textbf{Retrieval:} ChromaDB, cosine similarity, k=3 documents
    \item \textbf{Runs:} 5 independent trials (75 responses per approach)
    \item \textbf{Metrics:} ROUGE-L (primary), ROUGE-1/2, BLEU, METEOR
    \item \textbf{Statistics:} 95\% CI, paired t-tests, Cohen's d
\end{itemize}

\section{Results}

\subsection{Main Performance Comparison}

Table \ref{tab:main_results} presents results across all metrics. Contrary to our hypothesis, baseline (no retrieval) achieved the highest ROUGE-L score.

\begin{table}[h]
\centering
\small
\begin{tabular}{lccc}
\toprule
\textbf{Approach} & \textbf{R-L} & \textbf{BLEU} & \textbf{MET} \\
\midrule
\textbf{Baseline} & \textbf{.231} & .054 & \textbf{.310} \\
AI-Generated & .224 & \textbf{.057} & .290 \\
Human-Curated & .220 & .045 & .296 \\
Raw Sources & .208 & .039 & .275 \\
\bottomrule
\end{tabular}
\caption{Performance comparison (mean, 5 runs). R-L=ROUGE-L, MET=METEOR. All values $\pm$ 0.004-0.010 (95\% CI).}
\label{tab:main_results}
\end{table}

\subsection{Statistical Significance}

Table \ref{tab:statistical} shows pairwise comparisons. AI-generated achieved 102\% of human-curated performance (p=0.184, ns). Both significantly outperformed raw sources (p<0.05).

\begin{table}[h]
\centering
\small
\begin{tabular}{lcccc}
\toprule
\textbf{Comparison} & \textbf{Diff} & \textbf{d} & \textbf{p} & \textbf{Sig} \\
\midrule
AI vs Human & +2.1\% & 0.72 & .184 & No \\
AI vs Raw & +7.8\% & 3.19 & .002 & ** \\
Human vs Raw & +5.5\% & 1.89 & .014 & * \\
AI vs Baseline & -2.6\% & -1.07 & .074 & No \\
\bottomrule
\end{tabular}
\caption{Statistical comparisons (ROUGE-L). * p<.05, ** p<.01}
\label{tab:statistical}
\end{table}

\section{Discussion}

\subsection{Why Does Baseline Outperform RAG?}

We propose a \textbf{domain saturation hypothesis}: RAG benefit is inversely related to domain representation in LLM training data. For well-documented topics like smoking cessation:

\begin{enumerate}
    \item Baseline LLM knowledge is comprehensive from pre-training
    \item Retrieved context introduces noise rather than signal
    \item Context window constraints force omission of parametric knowledge
\end{enumerate}

This shifts the research question from ``Does RAG help?'' to ``\textit{When} does RAG help?''—a more nuanced and practically valuable inquiry.

\subsection{AI-Generated = Human-Curated}

The comparable performance (102\%, p=0.184) has practical implications: AI-generated knowledge bases can be produced at lower cost while maintaining quality. However, this is specific to smoking cessation and may not generalize to specialized domains requiring expert knowledge.

\subsection{Methodological Implications}

Our data leakage discovery demonstrates systemic risk in NLP evaluation. The contaminated test set produced completely reversed conclusions. We recommend:

\begin{itemize}
    \item Mandatory fuzzy string matching (<50\% threshold)
    \item Manual inspection of high-similarity cases
    \item Reporting similarity distributions in publications
\end{itemize}

Recent work has identified similar issues in popular benchmarks \cite{lewis2021questionanswering, dodge2021documenting}, suggesting this is a widespread problem.

\section{Limitations and Future Work}

\textbf{Limitations:}
\begin{itemize}
    \item Single domain (smoking cessation)
    \item Single LLM (GPT-4o-mini)
    \item Small test set (15 questions)
    \item Lexical metrics (not semantic accuracy)
    \item No human evaluation
\end{itemize}

\textbf{Future work:}
\begin{itemize}
    \item Test domain saturation hypothesis across domains with varying LLM knowledge coverage (rare diseases, emerging treatments)
    \item Human preference studies with smoking cessation counselors
    \item Larger test sets with increased statistical power
    \item Semantic metrics and factual accuracy evaluation
\end{itemize}

\section{Conclusion}

This study makes two primary contributions. \textbf{Methodologically}, we demonstrate that data leakage can fundamentally reverse research conclusions, serving as a critical warning for NLP evaluation. \textbf{Empirically}, we find that RAG benefit is domain-dependent: for well-documented topics, LLM parametric knowledge may suffice, and RAG overhead may not be justified.

When RAG is employed, AI-generated knowledge bases can match human-curated ones at reduced cost. These findings have implications for both researchers designing evaluations and practitioners deploying RAG systems in production.

Our materials (test set, results, code) are publicly available to support reproducibility and future research.

\section*{Acknowledgments}

We thank domain experts who validated test questions and the professor for guidance throughout this research.

\bibliography{custom}
\bibliographystyle{acl_natbib}

\appendix

\section{Complete Metrics}

Full ROUGE-1, ROUGE-2 results and per-question breakdowns are available in supplementary materials at: \url{https://github.com/[username]/rag-smoking-cessation-eval}

\section{Example Test Questions}

\begin{enumerate}
    \item ``My doctor says my lungs are already damaged from decades of cigarettes. What's the point of stopping now?''
    \item ``After avoiding tobacco for a full calendar year, which organs show the most significant recovery?''
    \item ``From a neuroscience perspective, explain the chemical chain reaction that occurs in the brain within seconds of nicotine inhalation.''
\end{enumerate}

\end{document}
